%% ============================================================
%% BENCHMARK CONSTRUCTION DETAILS
%% ============================================================

\section{Guideline Generation from Medical Literature}
\label{app:guideline_gen}

For each clinical concept in the dependency graph, we generate a structured natural-language specification by querying PubMed. Given a concept name (e.g., ``acute kidney injury''), we search PubMed for relevant clinical practice guideline articles and systematic reviews, retrieve 3--5 top-ranked results, and extract their abstracts and available full text. An LLM then synthesizes the retrieved literature into a structured specification containing:
\begin{enumerate}[label=(\roman*),nosep]
    \item the clinical definition and diagnostic criteria with explicit measurement thresholds (e.g., serum creatinine increase $\geq 0.3$~mg/dL within 48 hours per KDIGO guidelines),
    \item the required EHR measurements mapped to MIMIC-IV table names and item identifiers,
    \item severity staging criteria where applicable, and
    \item temporal windows for event detection.
\end{enumerate}
This pipeline produces guideline specifications for $N$ concepts spanning the measurement, scoring, and condition layers of the dependency graph.

\section{Clinical Knowledge Graph Construction}
\label{app:knowledge_graph}

We construct a directed acyclic dependency graph $\mathcal{G} = (V, E)$ where an edge $(v_i, v_j) \in E$ indicates that concept $v_j$ requires the output of concept $v_i$. The graph is built in three phases.

\textbf{Phase 1: Infrastructure nodes.} We parse SQL definitions from the MIMIC-Code repository, extracting references to upstream derived tables, raw MIMIC-IV tables, and \texttt{itemid} values via regex matching to produce infrastructure nodes (measurements, medications, severity scores, treatments) with exact computational dependencies.

\textbf{Phase 2: Clinical conditions.} We expand the graph by mining frequent ICD-10 diagnosis codes among MIMIC-IV ICU patients, filtering for conditions detectable from structured EHR data, and generating guideline specifications using the pipeline in Appendix~\ref{app:guideline_gen}. Dependencies for each condition are inferred by keyword-matching the specification text against infrastructure node descriptions and entity labels.

\textbf{Phase 3: UMLS expansion.} We map existing nodes to UMLS Concept Unique Identifiers and traverse one-hop Metathesaurus relations, filtering candidates by semantic type and ranking by centrality (the number of seed nodes connecting to each candidate).

The resulting graph contains multiple quality tiers—from infrastructure nodes with exact SQL-parsed dependencies, through conditions with guideline-validated dependencies, to automatically discovered conditions with heuristic dependency assignments—enabling both precise retrieval for core clinical concepts and scalability evaluation across a broad condition library.

%% ============================================================
%% ALGORITHM DETAILS
%% ============================================================

\section{Autoformalization Algorithm}
\label{app:algorithm}

\begin{algorithm}[t]
\caption{Autoformalize}
\label{alg:autoformalize}
\begin{algorithmic}[1]
\Require QA dataset $\mathcal{Q}$, database $\mathcal{C}$, guidelines $\mathcal{G}$, LLM $\mathcal{M}$, threshold $\tau$, max iterations $T$
\State $\mathcal{S} \leftarrow \textsc{CodeSession}(\mathcal{C}, \mathcal{G})$; \;
       $f^* \leftarrow \emptyset$; \; $\alpha^* \leftarrow 0$; \; $\text{mem} \leftarrow []$
\State $\mathit{msgs} \leftarrow [\textsc{SysPrompt}, \textsc{UserPrompt}(\mathcal{Q})]$
\For{$t = 1, \dots, T$}
    \If{context near limit}
        \State $\mathit{msgs} \leftarrow \textsc{Compress}(\text{mem} \cup \textsc{Summarize}(\mathit{msgs}))$
    \EndIf
    \State $r \leftarrow \mathcal{M}(\mathit{msgs})$; \; $\mathit{msgs} \mathrel{+}= [r]$
    \State $\mathcal{B} \leftarrow \textsc{ExtractCodeBlocks}(r)$; \; \textbf{if} $\mathcal{B}=\emptyset$ \textbf{continue}
    \State $\mathit{msgs} \mathrel{+}= [\textsc{Execute}(\mathcal{B}, \mathcal{S})]$
    \If{$\texttt{FINAL\_FUNCTION} \notin \mathit{msgs}[-1]$} \textbf{continue} \EndIf
    \State $f \leftarrow \mathcal{S}.\text{namespace}[\texttt{FINAL\_FUNCTION}]$
    \State $\alpha \leftarrow \textsc{QAEval}(f, \mathcal{Q}, \mathcal{C}, \mathcal{M})$
    \If{$\alpha > \alpha^*$} \; $f^* \leftarrow f$; \; $\alpha^* \leftarrow \alpha$ \EndIf
    \If{$\alpha \geq \tau$} \textbf{break} \EndIf
    \State $\mathit{msgs} \mathrel{+}= [\textsc{Feedback}(\alpha, \tau)]$
\EndFor
\State \Return $f^*$
\end{algorithmic}
\end{algorithm}

\paragraph{Context Compression.}
Because the ReACT loop may run for up to $T_{\max}$ iterations (default 100) with substantial code and query outputs at each step, the conversation history can exceed the LLM's context window. We address this with a rolling memory mechanism: when the estimated token count exceeds 25{,}000 tokens, the accumulated conversation is summarized by a separate LLM call into a compact memory document retaining key schema facts, the latest function source code, and a record of test results and error fixes. The full conversation history is then replaced by a two-message prompt (system instructions + memory resume). This compression is lossy but preserves the information most critical for continued refinement.

%% ============================================================
%% PROMPTS
%% ============================================================

\section{Prompts}
\label{app:prompts}

\subsection{Autoformalization Agent}
\begin{figure*}[h]
    \begin{AIbox}{Autoformalization System Prompt}

\begin{quote}
\ttfamily
\tiny
\vspace{2mm}

You are an expert clinical informaticist and Python programmer.

You can execute Python code by placing it inside <code> ... </code> XML tags.
The code runs in a persistent session with these pre-loaded:
  - `query\_db(sql)` — runs a SQL query on the DuckDB database, returns a
    pandas DataFrame
  - `pandas` (pd), `numpy` (np), `duckdb`, `datetime` are pre-imported
  - `DB\_PATH` — path to the DuckDB database


IMPORTANT DATABASE NOTES:
  
  \{Database Specific Notes\}

YOUR TASK:
Write a Python function that extracts critical information for the clinical concept described
in the guideline you are given.  The function should use `query\_db()` to
look up patient data from the database and return useful clinical information.

Steps:

1. **Explore the database schema** — list tables, inspect columns, look up
   dictionary/lookup tables (d\_labitems, d\_items, d\_icd\_diagnoses, etc.).
   Run sample queries to understand data formats.

2. **Write a Python function** — implement the guideline's logic.  The
   function should accept patient identifiers (e.g. subject\_id, hadm\_id,
   stay\_id) and return information that captures the clinical concept.
   The output format is flexible — return whatever best represents the
   concept (e.g., a value, dict, DataFrame, boolean, list, etc.).

3. **Test thoroughly** — call your function with several sample patients
   and verify the output looks reasonable.  Debug and fix any issues.
   Do NOT assign FINAL\_FUNCTION until you are confident it works.

4. **Save it** — once you are satisfied, write a **single, self-contained
   code block** that includes ALL imports, helper functions, and the main
   function definition, then assigns FINAL\_FUNCTION at the end.  This
   block must be runnable on its own — do NOT include test calls, print
   statements, or assertions in it.  
   
   Example:

       import pandas as pd

       def helper(x):
           ...

       def compute\_concept(subject\_id, hadm\_id):
           ...

           return value

       FINAL\_FUNCTION = compute\_concept

5. After you receive evaluation feedback, **revise** your function,
   test the revision, then submit a new self-contained code block
   assigning FINAL\_FUNCTION again.

IMPORTANT RULES:

- Your function MUST use `query\_db()` to access the database.


- Do NOT call `duckdb.connect()` directly — it will cause connection errors.


- Your function should accept patient identifiers and return information
  that usefully captures the clinical concept. Any output format is fine.
  
- The function must work for any valid patient, not just the test cases.


- Do NOT guess column names or item IDs — always look them up first.

- The code block that assigns FINAL\_FUNCTION must be entirely
  self-contained (all needed imports, helpers, and the function itself)
  and must NOT contain any test calls, print statements, or assertions.

Be methodical.  Explore first, code second.  
You may include multiple <code> blocks in a single response.

\end{quote}

\end{AIbox}
    \caption{The system prompt provided to the LLM agent during the autoformalization react loop. The prompt can be customized to include database specific information, such as important column names, primary key, etc.}
    \label{prompts:system_prompt}
    
\end{figure*}
\begin{figure*}[h]
    \begin{AIbox}{Autoformalization Task Prompt}

\begin{quote}
\ttfamily
\tiny
\vspace{2mm}

YOUR TASK --- build a function for the clinical concept `\{concept\_name\}':

You need to write a Python function that can help answer clinical
questions about `\{concept\_name\}' for individual patients.

Here are example questions your function will need to support:
\{sample\_questions\}

APPROACH:
1. **Search for guidelines** --- call `search\_guidelines()' to see what clinical
   guidelines are available, then use relevant keywords (e.g.
   `search\_guidelines("\{concept\_name\}")') to find the most relevant ones.
   Read each relevant guideline with `get\_guideline(name)' --- these contain
   clinical definitions, scoring criteria, and implementation details that
   are essential for getting the logic right.
2. **Check for reusable functions** --- call `search\_functions()' to see if any
   previously implemented functions could serve as helpers. Use
   `get\_function\_info(name)' to inspect their signatures and docstrings,
   and `load\_function(name)' to import ones you want to reuse.
3. **Explore the database schema** --- list tables, inspect columns, look up
   dictionary/lookup tables (d\_labitems, d\_items, d\_icd\_diagnoses, etc.).
   Run sample queries to understand data formats.
4. **Write a Python function** --- implement the clinical logic based on what
   you learned from the guidelines and the database schema.  The function
   should accept patient identifiers (e.g. subject\_id, hadm\_id, stay\_id)
   and return information that captures the clinical concept.
   The output format is flexible --- return whatever best represents the
   concept (e.g., a value, dict, DataFrame, boolean, list, etc.).
5. **Test thoroughly** --- call your function with several sample patients
   and verify the output looks reasonable.  Debug and fix any issues.
   Do NOT assign FINAL\_FUNCTION until you are confident it works.
6. **Save it** --- once satisfied, write a single self-contained code block
   with ALL imports, helpers, function def, and FINAL\_FUNCTION assignment.
   No test calls or print statements in that final block.

FINAL FUNCTION REQUIREMENTS:
  - It should return useful clinical information that captures the concept.
  - The input should involve some patient identifier (e.g. stay\_id)
  - The output can be any format (scalar, dict, DataFrame, etc.) --- choose
    whatever best represents the clinical concept.
  - Use `query\_db(sql)' or `query\_fhir(resource\_type, params)' to look up data.
  - Have a detailed docstring which explains the expected input and output.

Start by searching for guidelines relevant to `\{concept\_name\}'.

\end{quote}

\end{AIbox}
    \caption{The task prompt provided to the LLM agent to initiate autoformalization of a clinical concept. Template variables \{concept\_name\} and \{sample\_questions\} are filled at runtime.}
    \label{prompts:react_user_prompt}

\end{figure*}


\subsection{Iterative Feedback}
\begin{figure*}[h]
    \begin{AIbox}{QA Accuracy Feedback Prompt}

\begin{quote}
\ttfamily
\tiny
\vspace{2mm}

QA ACCURACY: \{qa\_accuracy\} (need >\{match\_threshold\})

\{details\}

Your function is evaluated by how well its output helps answer clinical
questions about suspected infection. Revise your function to produce more
clinically useful and informative output.

Test your revised function first, then submit a new self-contained code
block (all imports, helpers, function def) assigning FINAL\_FUNCTION.
Do NOT include test calls or print statements in that block.

\end{quote}

\end{AIbox}
    \caption{Feedback prompt sent to the agent when QA accuracy falls below the required threshold. Template variables \{qa\_accuracy\}, \{match\_threshold\}, and \{details\} are filled with the actual evaluation results and per-question failure details.}
    \label{prompts:comparison_feedback}

\end{figure*}

\begin{figure*}[h]
    \begin{AIbox}{Error Feedback Prompt}

\begin{quote}
\ttfamily
\tiny
\vspace{2mm}

Evaluating your function raised an error:
\{traceback\}

Fix the function, test it, then submit a new self-contained code block
assigning FINAL\_FUNCTION (with all imports and helpers, no test calls).

\end{quote}

\end{AIbox}
    \caption{Feedback prompt sent to the agent when the submitted function raises a runtime error during evaluation. The full Python traceback is included in \{traceback\}.}
    \label{prompts:error_feedback}

\end{figure*}

\begin{figure*}[h]
    \begin{AIbox}{Structural Feedback Prompts}

\begin{quote}
\ttfamily
\tiny
\vspace{2mm}

\textbf{(a) No code blocks detected:}

No <code> blocks found. Please write code inside <code>...</code> tags to explore the schema and build your function.

\vspace{3mm}

\textbf{(b) No FINAL\_FUNCTION assignment:}

You haven't set FINAL\_FUNCTION yet.
Once you have tested your function and are confident it works, write a
single self-contained code block with all imports, helpers, the function
definition, and `FINAL\_FUNCTION = <your\_func>' at the end.
Do NOT include test calls or print statements in that block.

\vspace{3mm}

\textbf{(c) Final iteration notice:}

IMPORTANT: This is your FINAL iteration. You MUST produce your best
FINAL\_FUNCTION now. Write a single self-contained code block containing
all imports, helpers, the function definition, and
`FINAL\_FUNCTION = <your\_func>' at the end.
Do NOT include test calls or print statements in that block.
Even if you are not fully confident, submit your best attempt now.

\end{quote}

\end{AIbox}
    \caption{Structural feedback prompts used to enforce agent output format. (a) is sent when the agent responds with no \texttt{<code>} blocks, (b) when the agent has not yet submitted a \texttt{FINAL\_FUNCTION}, and (c) when the maximum iteration budget is reached.}
    \label{prompts:structural_feedback}

\end{figure*}


\subsection{QA Evaluation}
\begin{figure*}[h]
    \begin{AIbox}{QA Evaluation System Prompt}

\begin{quote}
\ttfamily
\tiny
\vspace{2mm}

You are a clinical expert evaluating patient data for clinical concepts.

Execute Python code inside <code> ... </code> tags. The session pre-loads:
  - `query\_db(sql)' --- SQL query $\rightarrow$ pandas DataFrame
  - The concept function described in your task
  - `pandas' (pd), `numpy' (np), `datetime' pre-imported

Database: DuckDB with schemas `mimiciv\_hosp', `mimiciv\_icu', `mimiciv\_derived'.
Always use fully-qualified table names (e.g. `mimiciv\_hosp.admissions').
Never call `duckdb.connect()' --- use `query\_db(sql)'.
IDs: subject\_id (patient) $\rightarrow$ hadm\_id (admission) $\rightarrow$ stay\_id (ICU stay).

\end{quote}

\end{AIbox}
    \caption{System prompt for the QA evaluation agent. The agent uses the autoformalized concept function loaded into the session to answer clinical questions about individual patients.}
    \label{prompts:qa_system_prompt}

\end{figure*}

\begin{figure*}[h]
    \begin{AIbox}{QA Evaluation User Prompt}

\begin{quote}
\ttfamily
\tiny
\vspace{2mm}

Session tools:
- `query\_db(sql)' --- duckdb SQL query $\rightarrow$ pandas DataFrame
- `search\_functions(keyword="")' --- list available concept functions by name
- `load\_function(name)' --- load a concept function into the session
- `get\_function\_info(name)' --- view a function's signature and docstring
- `search\_guidelines(keyword="")' / `get\_guideline(name)' --- clinical guidelines

Concept functions relevant to this task are available --- use `search\_functions()' to discover them and `load\_function(name)' to load what you need. You should always check first to see if there are any useful functions for you to use.
\{func\_info\}
Patient: subject\_id=\{subject\_id\}, hadm\_id=\{hadm\_id\}, stay\_id=\{stay\_id\}

\{qa\_question\}

Use the available functions and/or `query\_db' to answer, then reply with: \{verdict\_format\}
For numerical questions, include as much decimal precision as possible.

\end{quote}

\end{AIbox}
    \caption{User prompt for the QA evaluation agent. Template variables are filled at runtime: \{func\_info\} provides signatures of relevant concept functions, \{subject\_id\}/\{hadm\_id\}/\{stay\_id\} identify the target patient, \{qa\_question\} is the clinical question, and \{verdict\_format\} specifies the answer format.}
    \label{prompts:qa_user_prompt}

\end{figure*}


%% ============================================================
%% Training Parameters
%% ============================================================

\section{Training details}
\label{app:train_details}

\begin{table}[t]
\centering
\small
\begin{tabular}{ll}
\toprule
\textbf{Category} & \textbf{Hyperparameter} \\
\midrule

\multicolumn{2}{c}{\textit{Model}} \\
Max prompt length & 1024 \\
Max response length & 16384 \\
Gradient checkpointing & Enabled \\
Model dtype & bfloat16 \\

\midrule
\multicolumn{2}{c}{\textit{Training}} \\
Algorithm & GRPO \\
Total epochs & 1 \\
Train batch size & 8 \\
PPO mini-batch size & 4 \\
Actor learning rate & $1 \times 10^{-5}$ \\
Clip ratio (low) & 0.20 \\
Clip ratio (high) & 0.28 \\
Clip ratio constant & 10.0 \\
Use KL in reward & False \\
KL coefficient & 0.0 \\
Use KL loss & False \\
KL loss coefficient & 0.0 \\
Dynamic batch size & Enabled \\

\midrule
\multicolumn{2}{c}{\textit{Rollout / Inference}} \\
Backend & vLLM (async) \\
Tensor parallel size & 2 \\
Responses per prompt (train) & 4 \\
Responses per prompt (val) & 1 \\
Temperature (val) & 1.0 \\
Top-p (val) & 0.6 \\
Max user turns & 5 \\
Max assistant turns & 5 \\
Multi-turn format & Hermes \\

\bottomrule
\end{tabular}
\caption{Training and inference hyperparameters for the MIMIC-IV agent using GRPO.}
\label{tab:hyperparams}
\end{table}

%% ============================================================
%% COMPUTE RESOURCES
%% ============================================================

\section{Compute Resources}
\label{app:compute}

All inference experiments were run on a \texttt{8xB200 Node}. Inference was served with \texttt{vllm}. Approximate runtimes per experimental condition are as follows:

\begin{itemize}[nosep]
    \item \textbf{Autoformalization (per concept):} 20 GPU-hours on average for a \texttt{Qwen3.5-27B} model run with efficient parallelization via topological ordering of concepts.
    \item \textbf{GRPO Finetuning:} \textcolor{red}{[TODO: Z]} GPU-hours total, ran on a subset of the training split (5 questions per concept) for a total of \textcolor{red}{[TODO: Z]} questions.
    \item \textbf{Zero-shot inference evaluation (56{,}700 QA pairs, one model):} 114 GPU-hours when using \texttt{Qwen3.5-27B} and parallel processing of 4 concepts at a time. Further gain could be made
\end{itemize}

Total compute across all reported experiments (all models, all baselines) is estimated at \textcolor{red}{[TODO: total]} GPU-hours. No substantial additional compute was used for preliminary or failed experiments beyond the reported runs.
